\documentclass[12pt]{article}

% Encoding and fonts
\usepackage[T1]{fontenc}
\usepackage[utf8]{inputenc}
\usepackage{lmodern}

% Page layout
\usepackage[margin=0.75in]{geometry}
\usepackage{multicol}
\usepackage{indentfirst}
\usepackage{setspace}

\setstretch{1.2}

% Mathematics
\usepackage{amsmath,amssymb,amsthm}

% Graphics and tables
\usepackage{graphicx}
\usepackage{booktabs}
\usepackage{tabularx}
\usepackage{array}
\usepackage{caption}

% Lists
\usepackage{enumitem}

% Colors and hyperlinks
\usepackage{xcolor}
\usepackage[colorlinks=true, urlcolor=blue]{hyperref}

% Paragraph formatting
\setlength{\parskip}{0.75em}

% Title Formatting
\title{Active Site Predictor Proposal}
\author{Duncan Ritchie Ian Carr}
\date{\today}

\makeatletter
\newcommand{\mytitle} {
	\begin{center}
		\LARGE\bfseries\@title
	\end{center}

	\begin{center}
		\begin{tabularx}{0.80\textwidth}{@{}L{1} C{1.2} C{1.2} R{0.6}@{}}
			\@author & CS 4800-102 & \@date & Waters
		\end{tabularx}
	\end{center}
}
\makeatother

% user defined functions
\newcolumntype{L}[1]{>{\hsize=#1\hsize\raggedright\arraybackslash}X}
\newcolumntype{C}[1]{>{\hsize=#1\hsize\centering\arraybackslash}X}
\newcolumntype{R}[1]{>{\hsize=#1\hsize\raggedleft\arraybackslash}X}

% formatting checks
% \overfullrule=1pt

\begin{document}

	\mytitle

	\section*{Background}

		\subsection*{Overview}

		The Central Dogma of Biology is the theory that information is encoded as
		DNA, transcribed as RNA, and translated into proteins. DNA is the basic
		building block of all living organisms. DNA or deoxyribose nucleic acids 
		are a class of molecules that is comprised (for brevity) four so-called 
		'bases'. These bases are the molecules adenine, cytosine, guanine, and 
		thymine, each represented by the single letter codes A, C, G, and T,
		respectively. DNA can then be represented by a string of these single 
		letter codes. 

		Some sections of DNA encode for genes that will later be translated into
		3D structures called proteins. However, to get to that point, DNA must 
		first be transcribed into RNA or ribonucleic acid. During transcription,
		the transcribed DNA sequence remains relatively unchanged, but in this
		chemical form, the RNA molecule, a sequence that retains the same
		information as the DNA sequence, is able to leave the organelle where
		all of a cell's DNA is contained.

		An RNA molecule is then able to travel into a cell's intracelluar matrix,
		or cytoskeleton, where it then undergos translation by a ribosome. The
		RNA molecule that contains the same information as the original DNA gene
		sequence adheres to a ribosome. The ribosome views a set of three bases
		along the RNA molecule and translates that set into a single amino acid.
		Once a set of three bases, or codon, is read and translated, the ribosome
		shifts three bases and translates the next codon into another amino 
		acid. The ribosome moves along the RNA molecule in this fashion until
		the end of the sequence is reached. The form of the translated RNA molecule
		is now an amino acid chain linked by peptide bonds, this is also called a
		polypeptide chain.

		Polypeptide chains are what proteins are made from, and each of the 20 
		amino acids that may be contained within any given chain are separated
		into their own classes based on the chemical properties they exhibit.
		These properties are the basis of protein folding. The polypeptide chain
		may start to fold into its final protein form, but will typically be
		transported to an additional organelle, the Golgi apparatus. In the Golgi
		apparatus, the polypeptide chain is correctly folded into its final form
		of a functional, or non-functional protein. These final proteins are
		the cell's workers, performing countless essential tasks your body
		needs to function normally, without issue. Functional proteins, or
		enzymes, are those that carry out some chemical process, whether that
		process be anabolic (creating one molecule from multiple substrate 
		subunits) or catabolic (breaking one substrate down into multiple).

		Molecules that bind to an enzyme are called substrates, and these
		so-called substrates can only bind to specific regions of the enzyme's
		3D structure. The regions where substrates may bind are called active
		sites, or catalytic centers. Substrates are able to bind to the active
		site of an enzyme because of the specific amino acids, or residues,
		placed along the surface of the protein and/or pocket of the active site.
		Proteins' active site are typically confirmed using technical lab
		techniques, but with this age of computing, predicting structures,
		properties, substrates, and biochemical pathways have reduced in-lab
		analysis. 

		One way to predict catalytic centers, or active sites, is to use a 
		motif detection approach. A motif is a well-known set of neighboring
		residues, or amino acids, that expresses some sort of attractive force
		for incoming substrates. Additionally, conservation between proteins
		in the same protein families has become an increasingly common approach
		to identify and predict how a protein may be used within a biochemical
		pathway.

		Finally, this project's purpose is to predict active sites for given
		proteins using Haskell as an underlying computational workhorse to 
		predict active sites, delivering the computed data to a C++ protein viewer
		highlighting the top-most predicted active site candidates. Duncan 
		Ritchie and Ian Carr will be working this semester to create a working
		prototype of this project detailed below.

		\subsection*{Features}

		While there's really only one main feature the user can interact with,
		this space will be used to describe the interal process of how this
		tool will work.

		At the command line, as the executable is called, an argument to the
		executable will be the accession number (a unique ID) of a protein.
		This accession number will be searched for and downloaded from the 
		RCSB PDB database. The file will then be automatically stripped of 
		all metadata and the transformation will result in only the structural
		data of the protein. The structural data will contain the spatial
		coordinates of individual atoms, the residue (molecule) those atoms
		are associated with, the sequence number the residue is associated with,
		and other unique identifiers for individual atoms.

		An initial C++ (phase 1) layer will access the RCSB PDB using REST API
		calls paired with RCSB PDB's web API. In phase 1, the C++ script will
		download the respective protein file, and pass that data to a Haskell
		script (phase 2) that will handle all parsing. After phase 2, the parsed
		data will then be analyzed by other Haskell scripts (phase 3) to handle
		internal computations like residue mapping and motif detection. Phase 3
		will also handle active site scoring, used to predict candidate active 
		site regions. Finally, in phase 4, the computed results will be passed
		to another C++ script to handle protein visualization and graphics 
		rendering using Vulkan. Vulkan is notoriously low--level and discrete,
		buffers will need to be manually allocated and data will need to be 
		manually uploaded and passed to the GPU. The resulting protein model will
		be rendered to the screen with a custom made low--level cross platform 
		Vulkan forward rendered rasterization engine. The most high scoring 
		candidate regions will then be highlighted within the protein model.
		The screen will be interactable, and you will be able to pan and rotate
		around the protein body, select which parts to render (everything,
		active sites only, or non--active sites only). The renderer will
		feature basic phong and dot based shading. Steps will be taken to ensure
		that this program is cross compatible with Linux and Windows devices.


		\subsection*{Final Product}

		In its final form, this tool will be a command line tool that will take,
		as input, the accession number of the protein in question (in RCSB PDB
		format), and as output, display a model of the the protein with thea
		predicted candidate active site regions highlighted within the model. 

		Additionally, testing of this tool will be carried out by comparing the
		active sites of known, experimentally validated proteins.

	\section*{Skills}

		\subsection*{Possessed}

		Already acquired skills that will help with the completion of this project
		include, but are not limited to:

		\begin{itemize}[noitemsep]
			\item Most of the Biology theory behind the project is known by
				one of the group members.
			\item Proficiency in both C++ and Haskell languages.
			\item Proficiency in Haskell parsing techniques, and functional
				programming paradigm.
			\item Proficiency in graphics programming, though known libraries
				may be higher-level compared to Vulkan.
			\item Both group members are passionate about each portion of the
				project, (Duncan - biology theory, Ian - graphics).
		\end{itemize}

		\subsection*{To Learn}

		Skills necessary to learn for the completion of this project include,
		but are not limited to:
		
		\begin{itemize}[noitemsep]
			\item Understanding build and metabuild systems to seamlessly work
				across different operating systems.
			\item Understanding Haskell and C++ FFI bridge system.
			\item Understanding REST APIs and the RCSB PDB web API.
			\item Understanding the Vulkan graphics library to visualize protein
				models.
			\item Optimizing scoring algorithm, ensuring accuracy and scientific
				reliability and validity.
			\item Understanding specific motif patterns and conservation to make
				reasonable scoring system.
		\end{itemize}

	\section*{Challenges}

	One major challenge will be explaining the biology theory to classmates
	as they may check in, ask questions, and suggest ideas or solutions.
	Getting the devised system to work cohesively together may prove to be
	a major challenge as each partner is working on different operating systems,
	as well as leveraging aggregated programming languages. Finally, the Vulkan
	ecosystem is quite low-level and will require a major time commitment to
	get working. Project and architecture design will be the main starting point,
	ensuring a modular, easy--to--debug system. Finally, Duncan has been taking
	time from his CS degree to focus on his minors for the past several semesters
	and will need to get back into the groove of programming, especially in C++.

	\section*{Division of Labor}

	As for how the labor of the project will be broken up between members, 
	Duncan will take on the role of Chief Algorithmic Engineer, while Ian will
	spearhead as the Graphics Department's Executive Director. All fun aside,
	a more detailed list of objectives and tasks by each team member is outlined 
	below.

	\section*{Literature Review}

	Similar work for the Haskell programming side includes a statistical model
	predicting active sites (1), that include other resources of interest, such as
	catalytic motifs and a basis for residue scoring (2). A separate paper includes
	a technique for dimensional analysis for finding clefts and pockets within
	or around the surface of the protein (3).

	\begin{enumerate}[noitemsep]
		\item \href{https://pubmed.ncbi.nlm.nih.gov/20080507/}{Active site prediction using evolutionary and structural information}
	\item \href{https://www.sciencedirect.com/science/article/pii/S0022283602010367}{Analysis of Catalytic Residues in Enzyme Active Sites}
	\item \href{https://pmc.ncbi.nlm.nih.gov/articles/PMC2700099/}{Fpocket: An open source platform for ligand pocket detection}
	\end{enumerate}
	
	Additionally, for the graphics section of this progect, a sample section of
	code as a "Hello World'' equivalent is available on GitHub and will be used
	as a reference (1). Another section of sample code showing a basic Blinn 
	Phong shading model, while displaying the processed affect (2). Finally, a
	paper explains the math behind modern day perspective projection matrices (3).

	\begin{enumerate}[noitemsep]
		\item \href{https://gist.github.com/Overv/7ac07356037592a1121225172d7d78f2d}{HelloTriangle.cc}
		\item \href{https://www.shadertoy.com/view/3d3cWN}{Blinn Phong Model}
		\item \href{https://dl.acm.org/doi/pdf/10.1145/356744.356750}{Planar Geometric Projections and Viewing Transformations}
	\end{enumerate}

	\section*{Stretch Goals}

	A few stretch goals for this project are:

	\begin{itemize}[noitemsep]
		\item Implementing conservation scoring by fetching protein conservation
			data, parallel processing of amino acid sequence, and joining
			conservation, structural, and motif scores on individual residues.
		\item Implementing shadow mapping into the renderer so it is more visually
			appealing. Showing depth in an image is an important part of graphics
			programming
		\item Different ways to render the protein molecules. There will be the
			base approach which simply renders spheres, and possibly a complex
			one which renders helixes
		\item Support for Macs and Apple devices using the MoltenVK layer
		\item Antialiasing (MSAA or FXAA.Undecided)
	\end{itemize}

\end{document}
